% Relatório: comparação da orquestração de agentes entre versões do NarraTer.
% Compilar: cd docs/relatorios && pdflatex orquestracao-versoes.tex (duas vezes).
% Os números vêm de dados/ (gerado por bench/orquestracao/agregar.py); não edite os .tex de dados.
\documentclass[11pt,a4paper]{article}
\usepackage[utf8]{inputenc}
\usepackage[T1]{fontenc}
\usepackage[brazil]{babel}
\usepackage{lmodern}
\usepackage[margin=2.3cm]{geometry}
\usepackage{booktabs}
\usepackage{tabularx}
\usepackage{graphicx}
\usepackage{xcolor}
\usepackage{pgfplots}
\usepackage{pgfplotstable}
\usepackage{hyperref}
\usepackage{enumitem}
\usepackage{float}
\pgfplotsset{compat=1.18}
\setlist{itemsep=2pt}

\definecolor{va}{HTML}{8A8F98}
\definecolor{vb}{HTML}{4F7FD9}
\definecolor{vc}{HTML}{8F7CFF}
\definecolor{vd}{HTML}{2E9E6B}
\newcommand{\rot}[1]{\textsf{\small[#1]}}

\title{Orquestração de agentes no NarraTer:\\ comparação entre quatro versões}
\author{Gabriel Cortez \and Claude Code}
\date{15 de outubro de 2026}

\begin{document}
\maketitle

\begin{abstract}
\input{texto/resumo.tex}
\end{abstract}

\section{Objetivo}
\input{texto/objetivo.tex}

\section{Versões comparadas}
\input{texto/versoes.tex}

\section{Método}
\input{texto/metodo.tex}

\section{Resultados}
\input{texto/resultados.tex}

\section{Discussão}
\input{texto/discussao.tex}

\section{Ameaças à validade}
\input{texto/ameacas.tex}

\section{Como reproduzir}
\input{texto/reproduzir.tex}

\end{document}
